\subsection{平展代数中的范数与迹}

设 A 是 K 上（有限）次数为 n 的平展代数。按定义，此时存在 K 的一个扩张 L 与 A 到 L 中的互不相同的同态 \(u_{1}, \ldots, u_{n}\) ，具有下列性质。

\begin{enumerate}
\def\labelenumi{\alph{enumi})}
\item
  \(\mathbf{A}\) 到 \(\mathbf{L}\) 中的每个同态都等于某个 \(u_{i}\) （V，第 29 页，推论）
\item
  存在 L-代数同构 \(u: \mathbf{A}_{(L)} \to L^n\) 使得
\end{enumerate}

\[
u (1 \otimes x) = \left(u _ {1} (x), \dots , u _ {n} (x)\right) \quad \text { for   all } \quad x \in A.
\]

此外，K 的每个代数闭扩张 L 都具有上述性质（V，第 30 页，命题 2）。

在以下的讨论中我们固定 \(L, u_1, \ldots, u_n\) 。设 \(x \in A\) ；我们将证明公式

\[
\operatorname{Tr} _ {\mathrm{A} / \mathrm{K}} (x) \cdot 1 = \sum_ {i = 1} ^ {n} u _ {i} (x), \quad \mathrm{N} _ {\mathrm{A} / \mathrm{K}} (x) \cdot 1 = \prod_ {i = 1} ^ {n} u _ {i} (x).\tag{3}
\]

设 \(v\) 为乘以 \(1 \otimes x\) 的映射，作用在 \(\mathbf{A}_{(L)}\) 上；相对于 \(\mathbf{A}_{(L)}\) 的那组基（即 \(u^{-1}\) 下 \(L^n\) 的典范基的像），线性映射 \(v\) 的矩阵是对角形的，其对角元为 \(u_1(x), \ldots, u_n(x)\) 。由此我们得出

\[
\operatorname{Tr} _ {\mathrm{A} _ {(\mathrm{L})} / \mathrm{L}} (1 \otimes x). 1 = \sum_ {i = 1} ^ {n} u _ {i} (x), \text { whence } \operatorname{Tr} _ {\mathrm{A} / \mathrm{K}} (x). 1 = \sum_ {i = 1} ^ {n} u _ {i} (x) \text { by } (1);
\]

范数的情形类似处理。

再设 \((x_{1},\ldots ,x_{n})\) 是由 A 的元素组成的序列，设 U 为矩阵

\[
\left(u _ {i} \left(x _ {j}\right)\right) _ {1 \leqslant i, j \leqslant n}
\]

并令 \((t_{ij}) = {}^t\mathrm{U} \cdot U\) 。由第一个公式 (3)，我们有

\[
\operatorname{Tr} _ {\mathrm{A} / \mathrm{K}} \left(x _ {i} x _ {j}\right) \cdot 1 = \sum_ {k = 1} ^ {n} u _ {k} \left(x _ {i} x _ {j}\right) = \sum_ {k = 1} ^ {n} u _ {k} \left(x _ {i}\right) u _ {k} \left(x _ {j}\right) = t _ {i j};
\]

取行列式，我们得到

\[
\mathrm{D} _ {\mathrm{A} / \mathrm{K}} (x _ {1}, \dots , x _ {n}) \cdot 1 = [ \det u _ {i} (x _ {j}) ] ^ {2}.\tag{4}
\]

命题 1.——设 A 是 K 上有限次的交换代数。则下列条件等价：

\begin{enumerate}
\def\labelenumi{\alph{enumi})}
\item
  代数 \(\mathbf{A}\) 是平展的。
\item
  存在 \(\mathbf{A}\) 的一组基，其判别式非零。
\item
  对 A 中每个 \(x \neq 0\) ，存在 A 中的 \(y\) 使得 \(\operatorname{Tr}_{\mathbf{A} / \mathbf{K}}(xy) \neq 0\) .
\end{enumerate}

此外，当这些条件满足时，A 的任何基的判别式都非零。

我们要证明：当假设 A 平展时，A 关于任意一组基 \((x_{1},\ldots,x_{n})\) （A 在 K 上的基）的判别式非零；这特别地确立了蕴含关系 \(a)\Rightarrow b)\) 。借助上述记号，由 (4)，只需证明矩阵 U 可逆，或等价地，证明线性方程组

\[
\sum_ {i = 1} ^ {n} \lambda_ {i} u _ {i} (x _ {j}) = 0 \quad (\text { for } 1 \leqslant j \leqslant n)\tag{5}
\]

在 L 中仅有解 \(\lambda_{1} = \cdots = \lambda_{n} = 0\) 。现在关系式 (5) 蕴含 \(\sum_{i=1}^{n} \lambda_{i} u_{i}(x) = 0\) （对所有 \(x \in \mathbf{A}\) ），从而 \(\lambda_{i} = 0\) （对 \(1 \leqslant i \leqslant n\) ），这依据的是

同态的线性无关性定理（V，第 27 页，定理 1）。

\(b)\) 与 \(c)\) 的等价性是下述一般引理的推论：

引理 1.——设 V 是 K 上有限维向量空间，B 是 V × V 上的一个双线性型。设 \((v_{1}, \ldots, v_{n})\) 是 V 在 K 上的一组基，并设 \(\Delta = \det B(v_{i}, v_{j})\) 。则 \(\Delta \neq 0\) 当且仅当：对 V 中每个 \(x \neq 0\) ，存在 V 中的 y 使得 \(B(x, y) \neq 0\) .

我们有 \(\Delta \neq 0\) 当且仅当线性方程组

\[
\sum_ {i = 1} ^ {n} \lambda_ {i} \mathrm{B} (v _ {i}, v _ {j}) = 0 \quad (1 \leqslant j \leqslant n)
\]

在 K 中只有解 \(\lambda_1 = \cdots = \lambda_n = 0\) 。若令 \(x = \sum_{i=1}^{n} \lambda_i v_i\) ，则上述方程组等价于 B \((x, v_j) = 0\) （对 \(1 \leqslant j \leqslant n\) ），也等价于——由于 \((v_1, \ldots, v_n)\) 是 V 在 K 上的一组基——B \((x, y) = 0\) 对所有 \(y \in V\) 成立，从而引理得证。

我们来证明条件 c) 蕴含 A 是既约的。设 x 是 A 中的幂零元；对任意 \(y \in A\) ，元素 xy 是幂零的，于是向量空间 A 的自同态 \(L_{xy}\) 是幂零的。而下述引理表明 \(\operatorname{Tr}(xy) = 0\) 对所有 \(y \in A\) 成立，从而在假设 c) 下有 x = 0。

引理 2.——设 V 是 K 上有限维向量空间，u 是 V 的一个幂零自同态，则 \(\operatorname{Tr}(u)=0\) .

对每个整数 \(n \geqslant 0\) ，令 \(V_n\) 为 \(u^n\) 的像。由于 \(u\) 是幂零的，存在整数 \(r \geqslant 0\) 使得 \(V_0 = V\) 、\(V_r = 0\) ，且 \(V_i \neq V_{i+1}\) （对 \(0 \leqslant i \leqslant r-1\) ）。令 \(d_i\) 为 \(V_{i-1}\) 的维数（对 \(1 \leqslant i \leqslant r\) ）。存在一组基 \((x_1, \ldots, x_d)\) （\(V\) 的），使得向量 \(x_j\) （满足 \(d - d_i < j \leqslant d\) ）构成 \(V_{i-1}\) 的一组基（对 \(1 \leqslant i \leqslant r\) ）。我们有 \(u(V_{i-1}) \subset V_i\) ，因此 \(u\) 在基 \((x_1, \ldots, x_d)\) 下的矩阵的对角元全为零。于是 \(\operatorname{Tr}(u) = 0\) ，引理得证。

最后我们证明 \(b)\) 蕴含 \(a)\) 。设 \((x_1, \ldots, x_n)\) 是 \(A\) 在 \(K\) 上的一组基，使得 \(D_{A/K}(x_1, \ldots, x_n) \neq 0\) 。设 \(K'\) 是 \(K\) 的一个扩张，\(A'\) 为 \(K'\) -代数（由 \(A\) 经标量扩张得到），且 \(x_i' = 1 \otimes x_i\) （对 \(1 \leqslant i \leqslant n\) ）。由公式 (2)（V，第 47 页），我们有 \(D_{A'/K'}(x_1', \ldots, x_n') \neq 0\) 。把前面的结果应用于 \(A'\) ，便知 \(A'\) 是既约的，从而代数 \(A\) 是平展的（V，第 34 页，定理 4）。

推论.——设 E 是 K 的有限次扩张。E 可分当且仅当存在 E 中的 a 使得 \(\operatorname{Tr}_{E/K}(a) \neq 0\) .

该条件由命题 1 知是必要的。反之，假设存在 \(a \in E\) 使得 \(\operatorname{Tr}_{E/K}(a) \neq 0\) 。给定 \(x \neq 0\) （\(E\) 中的），若令 \(y = ax^{-1}\) ，则 \(\operatorname{Tr}_{E/K}(xy) \neq 0\) 。于是命题 1 表明 \(E\) 是 \(K\) 上的一个平展代数，从而是 \(K\) 的可分扩张。

